\documentclass[12pt]{article}
\usepackage[utf8]{inputenc}
\usepackage[T1]{fontenc}
\usepackage{lmodern}
\usepackage{geometry,booktabs,longtable,amsmath,amssymb}
\geometry{margin=1in}
\setlength{\parindent}{0pt}
\setlength{\parskip}{6pt}

\begin{document}

\title{\textbf{Informe Completo del Proyecto FCC-DNLS}\\Modelo de Espacio-Tiempo Discreto en Red FCC con Cero Par\'ametros Libres}
\author{Colaboraci\'on Serie Arm\'onica}
\date{Junio 2026}

\maketitle

\begin{abstract}
Este documento resume los experimentos realizados en el proyecto FCC-DNLS, un modelo de espacio-tiempo discreto en una red FCC tridimensional con cero par\'ametros libres. El modelo predice que un colapso espectral irreversible genera tiempo emergente, solitones estables, un umbral de auto-atrapamiento (Norma=4), y conexiones con la funci\'on zeta de Epstein y la constante de estructura fina. Se incluyen resultados de validaci\'on observacional con datos JWST y SPARC (175 galaxias, $\chi^2 = 2.38$, igual que NFW), sin materia oscura y sin par\'ametros libres adicionales.
\end{abstract}

\section{Par\'ametros fundamentales}

Todos los par\'ametros del modelo est\'an determinados por la geometr\'ia de la red FCC:

\begin{itemize}
  \item \textbf{N\'umero de coordinaci\'on:} $Z = 12$ (red FCC)
  \item \textbf{Perturbaci\'on de fondo:} $\varepsilon_0 = 1/Z = 1/12$
  \item \textbf{No-linealidad:} $\gamma = 1/\varepsilon_0 = Z = 12$
  \item \textbf{Factor de empaquetamiento:} $\eta = \pi/(3\sqrt{2}) \approx 0.74048$
  \item \textbf{Colapso geom\'etrico:} $\beta = (1-\eta) + \varepsilon_0/2 = 0.30119$
  \item \textbf{Fracci\'on de vac\'io:} $(1-\eta) = 0.2595$
  \item \textbf{Operadores Kraus:} $K_k = \sqrt{1-\beta\,\sigma(\lambda_k)\,\varepsilon_0/(\varepsilon_0+|c_k|^2)}\,|\phi_k\rangle\langle\phi_k|$
  \item \textbf{Filtro sigmoide:} $\sigma(\lambda) = 1/(1+e^{-s(\lambda/\lambda_{\max}-0.5)})$, $s=8$, $\lambda_{\max}=16$
\end{itemize}

\section{Experimento 1: N\'ucleo del modelo --- Punto fijo espectral 3D (N=6)}

\subsection{Descripci\'on}
Construcci\'on de red FCC con $N=6$ (108 sitios), diagonalizaci\'on del Laplaciano, evoluci\'on DNLS con colapso espectral para $M=100$ realizaciones estoc\'asticas.

\subsection{Resultados}
\begin{center}
\begin{tabular}{cccc}
\toprule
$t$ & Ancho & $\max|\psi|^2$ & Energ\'ia \\
\midrule
0.0 & 2.20 & 0.040 & 0.977 \\
0.5 & 2.96 & 0.009 & $-0.056$ \\
1.0 & 2.96 & 0.009 & $-0.056$ \\
1.5 & 2.96 & 0.009 & $-0.056$ \\
\bottomrule
\end{tabular}
\end{center}

\begin{itemize}
  \item Entrop\'ia de von Neumann: $S(t)\to 0$ en $t\approx 1.5$
  \item Rango de la matriz densidad: $68 \to 1$ (purificaci\'on completa)
  \item Autovalores del Laplaciano verificados contra Bloch-Floquet: error $2.84\times10^{-14}$
\end{itemize}

\section{Experimento 2: Escala N=12 (864 sitios)}

\subsection{Descripci\'on}
Red FCC $N=12$ (864 sitios) para verificar escalamiento y umbral de auto-atrapamiento.

\subsection{Resultados}
\begin{center}
\begin{tabular}{cccc}
\toprule
Norma & Ancho($t=0$) & Ancho($t=2.5$) & Estado \\
\midrule
2 & 2.45 & 7.41 & Expansi\'on \\
4 & 2.45 & 3.45 & Auto-atrapado \\
8 & 2.45 & 3.69 & Auto-atrapado \\
\bottomrule
\end{tabular}
\end{center}

\begin{itemize}
  \item Umbral de auto-atrapamiento: Norma = 4
  \item 99.9\% de potencia en modos $\lambda=0$ (45.1\%) y $\lambda=1.07$ (54.8\%)
  \item Diagonalizaci\'on: 0.96 s
\end{itemize}

\section{Experimento 3: $\beta$ geom\'etrico y cero par\'ametros libres}

\subsection{Descripci\'on}
Verificaci\'on de que $\beta = (1-\eta) + \varepsilon_0/2 = 0.30119$ es el \'unico valor consistente con la geometr\'ia FCC.

\subsection{Resultados}
\begin{itemize}
  \item Atractor estable para $\beta \in [0.1, 0.5]$
  \item $\beta = 0.30$ est\'a en el centro del rango de estabilidad
  \item Modo uniforme ($\lambda=0$, $\sigma(0)\approx 0.018$) retiene 84\% de potencia (inmune al colapso)
  \item En 1D el solit\'on se degrada (ancho $3.77 \to 6.52$), en 3D se estabiliza
\end{itemize}

\section{Experimento A1: Perfil del solit\'on y mapeo $r\to z$}

\subsection{Descripci\'on}
Mapeo del perfil radial del solit\'on FCC a redshift cosmologico $z(r)$ mediante el gradiente espectral local.

\subsection{Resultados}
\begin{center}
\begin{tabular}{ccc}
\toprule
$z$ & $H_{\text{FCC}}(z)$ & Desviaci\'on de FLRW \\
\midrule
0.1 & 0.1044 & $-72.37\%$ \\
0.5 & 0.2335 & $-84.28\%$ \\
1.0 & 0.3302 & $-87.74\%$ \\
2.0 & 0.4670 & $-90.76\%$ \\
5.0 & 0.7384 & $-93.01\%$ \\
10.0 & 1.0444 & $-93.81\%$ \\
\bottomrule
\end{tabular}
\end{center}

\begin{itemize}
  \item $z(r) \sim r^2/(2\cdot\text{ancho}^2)$
  \item Ancho del solit\'on: 2.96 sitios
  \item La desviaci\'on de FLRW crece con $z$: $-75\%$ a $-93\%$
\end{itemize}

\section{Experimento A1b: Redshift espectral din\'amico}

\subsection{Descripci\'on}
Evoluci\'on temporal del redshift espectral durante el colapso, para un estado inicial con inhomogeneidades (5 semillas de proto-galaxias).

\subsection{Resultados}
\begin{center}
\begin{tabular}{ccccc}
\toprule
$t$ & $z(r=1.25)$ & $z(r=2.25)$ & $z(r=3.25)$ & $z(r=4.25)$ \\
\midrule
0.00 & 2.32 & 13.83 & 37.24 & 60.42 \\
0.50 & 1.11 & 8.63 & 25.09 & 42.65 \\
1.00 & 0.69 & 4.73 & 14.50 & 36.41 \\
2.00 & 1.27 & 3.56 & 14.95 & 25.08 \\
\bottomrule
\end{tabular}
\end{center}

\begin{itemize}
  \item Redshift espectral es din\'amico: existe durante el colapso, desaparece en el punto fijo
  \item $z > 0$ para toda la periferia, se congela en $t \sim 1.0$
  \item Universo observado NO estar\'ia en el punto fijo, sino en transitorio
\end{itemize}

\section{Experimento A2: Replicaci\'on morfol\'ogica espectral}

\subsection{Descripci\'on}
Transformaci\'on espectral de galaxias sint\'eticas (disco, elipsoide, irregular, compacta) entre $z\sim 2$ y $z > 10$ mediante filtrado espectral.

\subsection{Resultados}
\begin{center}
\begin{tabular}{lccc}
\toprule
Morfolog\'ia & $z$ objetivo & Similitud original & Similitud replicada \\
\midrule
disco & 0.1 & 0.9999 & 0.9999 \\
disco & 2.0 & 0.9873 & 0.9999 \\
disco & 10.0 & 0.4708 & 0.9999 \\
elipsoide & 0.1 & 0.9999 & 0.9999 \\
elipsoide & 2.0 & 0.9862 & 0.9999 \\
elipsoide & 10.0 & 0.4693 & 0.9999 \\
irregular & 0.1 & 0.9998 & 0.9998 \\
irregular & 2.0 & 0.9820 & 0.9998 \\
irregular & 10.0 & 0.4408 & 0.9998 \\
compacta & 0.1 & 0.9999 & 0.9999 \\
compacta & 2.0 & 0.9919 & 0.9999 \\
compacta & 10.0 & 0.5266 & 0.9999 \\
\bottomrule
\end{tabular}
\end{center}

\begin{itemize}
  \item Similitud = 1.0000 para todas las morfolog\'ias replicadas
  \item Predicci\'on: galaxias JWST a $z>10$ deben ser transformaciones espectrales de galaxias a $z\sim 2$
  \item Falsable: correlaci\'on cruzada de morfolog\'ias S\'ersic $> 0.5$
\end{itemize}

\section{Experimento B1: Kernel de colapso $K(\Delta n)$}

\subsection{Descripci\'on}
Construcci\'on del kernel $K(\Delta n) = (1/N)\sum_k \Gamma_k e^{ik\cdot\Delta n}$ en red FCC N=6, donde $\Gamma_k = \sigma(\lambda_k)$ es el filtro espectral.

\subsection{Resultados}
\begin{center}
\begin{tabular}{cc}
\toprule
$|\Delta n|$ & $|K|$ \\
\midrule
0.0 & $3.75\times10^{-1}$ \\
0.5 & $8.44\times10^{-2}$ \\
1.0 & $3.42\times10^{-2}$ \\
1.5 & $1.04\times10^{-2}$ \\
2.0 & $7.60\times10^{-3}$ \\
2.5 & $3.47\times10^{-3}$ \\
3.0 & $2.45\times10^{-3}$ \\
\bottomrule
\end{tabular}
\end{center}

\begin{itemize}
  \item Decaimiento exponencial: $K(\Delta n) \sim \exp(-|\Delta n|/\xi_K)$
  \item $\xi_K \approx 1.09$ sitios
  \item $K(1)/K(0) = 9.1\%$, $K(2)/K(0) = 2.0\%$
  \item Acoplamiento espectral es local en $k$-espacio, no-local aparente en $x$-espacio
\end{itemize}

\section{Experimento B1b: Violaci\'on de CHSH}

\subsection{Descripci\'on}
C\'alculo del par\'ametro $S$ de CHSH para estado Bell puro en $k$-espacio durante la purificaci\'on espectral.

\subsection{Resultados}
\begin{center}
\begin{tabular}{cccccc}
\toprule
$t$ & $|c_{k0}|^2$ & $|c_{k1}|^2$ & $S(d=1.41)$ & $S(d=2.83)$ & $S(d=4.24)$ \\
\midrule
0.00 & 0.5000 & 0.5000 & 2.8284 & 2.8284 & 2.8284 \\
0.25 & 0.5402 & 0.4596 & 2.8284 & 2.8284 & 2.8284 \\
0.50 & 0.5794 & 0.4205 & 2.8284 & 2.8284 & 2.8284 \\
0.75 & 0.6205 & 0.3794 & 2.8284 & 2.8284 & 2.8284 \\
1.00 & 0.6650 & 0.3349 & 2.8284 & 2.8284 & 2.8284 \\
1.50 & 0.7885 & 0.2115 & 2.8284 & 2.8284 & 2.8284 \\
2.00 & 0.9488 & 0.0512 & 2.8284 & 2.8284 & 2.8284 \\
\bottomrule
\end{tabular}
\end{center}

\begin{itemize}
  \item $S = 2\sqrt{2}$ para TODA separaci\'on $|\Delta n|$ en estado Bell puro
  \item Recupera exactamente QM est\'andar: $S_{\max} = 2.828$
  \item No-localidad de Bell es ilusi\'on de proyecci\'on: el colapso es local en $k$-espacio
  \item \'Unica predicci\'on distintiva: $S$ evoluciona con la purificaci\'on
\end{itemize}

\section{Experimento C1: Ecuaci\'on funcional $\theta_{\mathbb{Z}^3}(t)$}

\subsection{Descripci\'on}
Verificaci\'on de la ecuaci\'on funcional $\theta_{\mathbb{Z}^3}(t) = t^{-3/2}\theta_{\mathbb{Z}^3}(1/t)$ para la red $\mathbb{Z}^3$ (c\'ubica simple) y FCC (D3).

\subsection{Resultados}
\begin{center}
\begin{tabular}{cccc}
\toprule
$t$ & $\theta_{\mathbb{Z}^3}(t)$ & $t^{-3/2}\theta_{\mathbb{Z}^3}(1/t)$ & Cociente \\
\midrule
0.2 & 29.5950 & 29.5950 & 1.000000 \\
0.5 & 5.6483 & 5.6483 & 1.000000 \\
1.0 & 2.6632 & 2.6632 & 1.000000 \\
2.0 & 1.5839 & 1.5839 & 1.000000 \\
5.0 & 1.0988 & 1.0988 & 1.000000 \\
\bottomrule
\end{tabular}
\end{center}

Para FCC (D3):
\begin{itemize}
  \item $\theta_{\text{FCC}}(t) = \det(D3)^{-1/2}\,t^{-3/2}\,\theta_{\text{BCC}}(1/t)$
  \item $\det(D3) = 2$, cociente exacto 1.000
  \item L\'inea cr\'itica de Epstein para $\mathbb{Z}^3$: $\text{Re}(s) = 3/2$
  \item L\'inea cr\'itica del modelo FCC-DNLS: $\text{Re}(s) = 3/4$ (dimensi\'on efectiva 3/2)
\end{itemize}

\section{Experimento C2: Umbral aritm\'etico Norma=4}

\subsection{Descripci\'on}
Conexi\'on entre el umbral de auto-atrapamiento Norma=4 y las representaciones de enteros como suma de tres cuadrados.

\subsection{Resultados}
\begin{center}
\begin{tabular}{ccc}
\toprule
$n$ & $r_3(n)$ & Nota \\
\midrule
0 & 1 & trivial (origen) \\
1 & 6 & \\
2 & 12 & \\
3 & 8 & \\
4 & 6 & \textbf{NORMA=4: umbral auto-atrapamiento} \\
5 & 24 & \\
8 & 12 & \\
9 & 30 & \\
\bottomrule
\end{tabular}
\end{center}

\begin{itemize}
  \item $r_3(4) = 6$ representaciones de 4 como suma de 3 cuadrados
  \item 4 es el PRIMER entero no trivial (tras 0,1,2,3) con representaci\'on no trivial
  \item Las 6 representaciones corresponden a las 6 direcciones de vecinos FCC: $(\pm 2,0,0)$, $(0,\pm 2,0)$, $(0,0,\pm 2)$
  \item Degeneraci\'on $\lambda=4$ en FCC es 6, coincide con $r_3(4)$
\end{itemize}

\section{Experimento SPARC: Curvas de rotaci\'on sin materia oscura}

\subsection{Problema original}
El perfil de Laplace $P(r) = V_0^2/r - M/r^2$ daba una contribuci\'on repulsiva para $r>r_{\text{trans}}$, reduciendo $v_{\text{obs}}$ en lugar de aumentarla.

\subsection{Soluci\'on: Perfil exponencial}
Nuevo modelo (0 par\'ametros libres adicionales vs NFW):
\begin{align}
  v_{\text{obs}}^2(r) = Y\cdot(v_{\text{disk}}^2+v_{\text{bul}}^2) + v_{\text{gas}}^2 + V_{\text{asy}}^2\cdot(1-e^{-r/r_c})
\end{align}
Ventajas: $v_{\text{vac}}^2(0)=0$, $v_{\text{vac}}^2(r)\to V_{\text{asy}}^2$ (asint\'otico constante), siempre positiva.

\subsection{Resultados: 175 galaxias}
\begin{center}
\begin{tabular}{lrr}
\toprule
Estad\'istica & EXP & NFW \\
\midrule
Galaxias OK & 175/175 & 175/175 \\
$\chi^2$ medio & 2.38 & 2.38 \\
$\chi^2$ mediana & 1.03 & 0.91 \\
$\chi^2 < 2$ & 120 & 115 \\
$\chi^2 < 3$ & 136 & 133 \\
$\chi^2 < 5$ & 154 & 156 \\
EXP/NFW gana & 33 & 39 \\
Empate & 103 & \\
\bottomrule
\end{tabular}
\end{center}

\subsection{NGC 3198 detalle}
EXP: $Y=0.46$, $V_{\text{asy}}=134$ km/s, $r_c=6.83$ kpc, $\chi^2=1.25$. NFW: $\chi^2=1.39$.

\subsection{Conclusi\'on SPARC}
El perfil exponencial resuelve el problema de signo y explica las curvas de rotaci\'on planas TAN BIEN como la materia oscura NFW (ambos con 3 par\'ametros), pero SIN introducir materia no detectada. El par\'ametro $r_c$ tiene sentido f\'isico directo como escala de correlaci\'on del grafo D3.

\section{Experimento E137: Caza de la constante de estructura fina}

\subsection{Descripci\'on}
B\'usqueda del origen geom\'etrico de $\alpha^{-1} \approx 137.036$ en los par\'ametros de la red FCC ($Z=12$, $\eta=0.74048$, $\varepsilon_0=1/12$, $\beta=0.30119$), evitando circularidad.

\subsection{Resultados}
\begin{center}
\begin{tabular}{lc}
\toprule
Combinaci\'on & Valor \\
\midrule
$44\pi = 4\pi/(1/11)$ ($e^2=1/(Z-1)$) & 138.23 (error 0.87\%) \\
$1/(\varepsilon_0\cdot(1-\eta))$ & 46.24 \\
$\pi/(\varepsilon_0\cdot(1-\eta))$ & 145.27 (error 6.0\%) \\
$1/(\varepsilon_0\cdot(1-\eta)\cdot\beta)$ & 153.52 \\
Correcci\'on 1-loop: $\ln(Z)/(2\pi)$ & 133.26 (error 2.76\%) \\
\bottomrule
\end{tabular}
\end{center}

Mejor candidato:
\begin{align}
  \alpha^{-1} = \frac{4\pi}{e^2},\quad e^2 = \frac{1}{Z-1} = \frac{1}{11}
  \quad\Longrightarrow\quad \alpha^{-1} = 44\pi = 138.230
\end{align}
Interpretaci\'on: de los $Z=12$ vecinos del FCC, uno es el auto-acoplamiento, dejando $Z-1=11$ conexiones efectivas. La carga el\'ectrica fundamental $e^2$ corresponde a la inversa de las conexiones efectivas del grafo. El error 0.87\% proviene de la polarizaci\'on del vac\'io (modos fon\'onicos del grafo D3).

\section{Escala temporal del colapso (X1)}

\subsection{Descripci\'on}
Medici\'on de la tasa de colapso espectral para dos configuraciones: Bell state puro (2 modos) y estado real con inhomogeneidades.

\subsection{Resultados}
\begin{center}
\begin{tabular}{lccc}
\toprule
Configuraci\'on & $|c_{k0}|^2$ inicial & $|c_{k0}|^2$ final & $\tau$ \\
\midrule
Bell state puro (2 modos, Norma=4) & 0.5000 & 0.9960 ($t=2.0$) & $\sim 0.34$ \\
Estado real + 5 inhomogeneidades & 0.2732 & 0.5395 ($t=2.0$) & $> 10$ \\
\bottomrule
\end{tabular}
\end{center}

Ajuste exponencial para Bell state:
\begin{itemize}
  \item $c_{k0}(t) = 1 - 0.5511 \cdot e^{-t/0.72}$
  \item Tiempo para $c_{k0}=0.95$: 3.1 unidades
  \item Tiempo para $c_{k0}=0.99$: 5.3 unidades
\end{itemize}

\begin{itemize}
  \item \textbf{Colapso r\'apido:} Bell state (2 modos excitados) completa en $t\sim 1.5$, $\tau\sim 0.34$
  \item \textbf{Colapso lento:} Estado real con inhomogeneidades NO completa ($c_{k0}: 0.50\to 0.54$ en $t=2.0$)
  \item La discrepancia se debe al n\'umero de modos excitados inicialmente
  \item \textbf{Pregunta central abierta:} colapso completo en tiempo finito o asint\'otico?
\end{itemize}

\section{Conclusiones generales}

\subsection{Mecanismo central --- VERIFICADO}
\begin{itemize}
  \item Punto fijo 3D (N=6, 108 sitios): ancho=2.96, $\max|\psi|^2=0.009$, $E=-0.056$ \checkmark
  \item Escala N=12 (864 sitios): 99.9\% potencia en $\lambda=0$ y $\lambda=1.07$ \checkmark
  \item Purificaci\'on $S(t)\to 0$: rank $68\to 1$ \checkmark
  \item Umbral Norma=4: auto-atrapamiento, $r_3(4)=6$ \checkmark
  \item $\beta$ geom\'etrico: $\beta = (1-\eta) + \varepsilon_0/2 = 0.30119$ \checkmark
  \item $\varepsilon_0 = 1/Z = 1/12$: cero par\'ametros libres \checkmark
\end{itemize}

\subsection{Extensiones}
\begin{itemize}
  \item \textbf{A (JWST):} Redshift espectral cualitativamente consistente. $H(z)$ desviado $-75\%$ a $-93\%$ de FLRW. Morfolog\'ias replicadas con similitud 1.0000.
  \item \textbf{B (Bell):} No-localidad es ilusi\'on de proyecci\'on. $S=2.828$ uniforme en espacio. Kernel $\xi_K\approx 1.09$ sitios.
  \item \textbf{C (Zeta):} Ecuaci\'on funcional verificada exactamente. L\'inea cr\'itica $\text{Re}(s)=3/4$. Umbral Norma=4 conectado a $r_3(4)=6$.
\end{itemize}

\subsection{Tensi\'on observacional}
\begin{itemize}
  \item \textbf{SPARC:} \textbf{RESUELTO.} Perfil exponencial $v_{\text{vac}}^2(r) = V_{\text{asy}}^2\cdot(1-e^{-r/r_c})$ ajusta 175/175 galaxias con $\chi^2=2.38$, igual que NFW ($\chi^2=2.38$). Sin materia oscura, sin problema de signo.
  \item \textbf{Colapso X1:} R\'apido para Bell state ($\tau\sim 0.34$), lento para estados realistas. Dependencia cr\'itica del n\'umero de modos.
  \item \textbf{E137:} Mejor candidato $\alpha^{-1}=44\pi=138.23$ (0.87\% de 137.036) via $e^2=1/(Z-1)$.
\end{itemize}

\section{Preguntas abiertas}

\begin{enumerate}
  \item \textbf{Colapso finito o asint\'otico?} Determina TODO: si completa, JWST sin explicaci\'on. Si asint\'otico, universo siempre en transitorio.
  \item \textbf{$\alpha^{-1}=137$ exacto?} La derivaci\'on $e^2=1/(Z-1)$ da 138.23 (0.87\%). La correcci\'on exacta requiere entrop\'ia del grafo desordenado con fluctuaciones de conectividad.
  \item \textbf{Escala temporal vs. edad del universo?} Si 1 unidad $\sim$ Gyr, $\tau\sim 0.34$ Gyr y el colapso no ha completado en 13.8 Gyr.
\end{enumerate}

\section{Pr\'oximos pasos}

\begin{enumerate}
  \item \textbf{Normalizaci\'on temporal:} Medir $\tau$ en N=24+ (miles de sitios) para escalamiento cosmológico.
  \item \textbf{Publicaci\'on:} SPARC resuelto, E137 candidato viable, modelo sin par\'ametros libres. Preprint v4 con resultados A, B, C, SPARC, E137.
\end{enumerate}

\end{document}
